%%%%%%%%%%%%%%%%%%%%%%%%%%%%%%%% BLOCK: tikhonov %%%%%%%%%%%%%%%%%%%%%%%%%%%%%%
\subsection*{Residual covariance and regularization}

Earlier versions of this work motivated Tikhonov regularization of the
residual covariance by rank deficiency: with $J$ features and $N$ training
texts, $J > N$ makes $\hat{\Sigma}$ singular and a ridge term
$\lambda \mathbf{I}$ necessary for invertibility.  That justification no
longer applies and is replaced here.

The power analysis of the present corpus (S2~Table) shows that a screening
threshold retaining 35 features at $N = 23$ carries an estimated false
discovery proportion above 0.5.  The defensible operating point is
$\alpha = 0.05$, retaining 14 features, with Benjamini-Hochberg at
$q = 0.10$ retaining 7 as a conservative band.  At $J = 14 < N = 23$ the
sample residual covariance is no longer rank-deficient, and regularization
is not required for invertibility.

Regularization is nonetheless retained, for a different and better reason.
The residual covariance is estimated from 23 observations and is
correspondingly noisy; its smallest eigenvalues are estimated worst, and
the likelihood of Eq.~\ref{eq:loglik} weights each direction by the inverse
of its eigenvalue.  Shrinking the spectrum toward its mean therefore
controls the influence of directions the data cannot resolve.  We select
$\lambda$ by inner leave-one-out on the training folds rather than fixing
it a priori.  The multiverse analysis reported under \nameref{sec:validation}
shows why this matters: without shrinkage, adding chance-correlated
features narrows credible intervals rather than widening them.

%%%%%%%%%%%%%%%%%%%%%%%%%%%% BLOCK: hbvi_results %%%%%%%%%%%%%%%%%%%%%%%%%%%%%%
\subsection*{Pentateuchal source dating}

Source-composite results from both model families are reported together
under \nameref{sec:results}, after the validation analysis that determines
how they should be read.  We defer them deliberately.  The empirical
results of this section in earlier versions of this work were computed
under a design in which each held-out unit was dated using its own
scholarly date as a prior; the resulting figures are not comparable to
those reported here and are not repeated.

%%%%%%%%%%%%%%%%%%%%%%%% BLOCK: prior_sensitivity_intro %%%%%%%%%%%%%%%%%%%%%%%
The prior sensitivity analysis that follows was originally conceived as a
robustness check.  In light of the leakage analysis under
\nameref{sec:validation} it carries more weight than that: the degree to
which a posterior is prior-dominated is not a peripheral diagnostic but the
central quantity governing whether a reported date is evidence about a text
or a restatement of an assumption.  We accordingly report the
classification for every out-of-sample unit rather than for a selected
subset.

%%%%%%%%%%%%%%%%%%%%%%%%%%%%% BLOCK: conclusion %%%%%%%%%%%%%%%%%%%%%%%%%%%%%%%
\section*{Conclusion}

Whether the morphosyntax of Biblical Hebrew carries a usable diachronic
signal has been disputed for two decades, with one side treating linguistic
dating as an established instrument and the other arguing that its apparent
successes are circular.  The analysis reported here supports a position
available to neither camp as usually stated: the signal is real, the
circularity is real, and the two facts are compatible.

The signal is real in a specific and testable sense.  Under a design in
which no text contributes to its own prediction and significance is
established by permuting date labels through the entire pipeline, Hebrew
morphosyntactic features order text pairs chronologically at 69.8\%
accuracy across 295 pairs, and leave-one-out rank correlations between true
and predicted date reach $+0.58$ and $+0.49$ in two independently
specified model families.  These are not large effects, but they are not
chance, and they replicate.

The circularity is real in an equally specific sense.  Holdout accuracies
of order 10--30~yr, reported in this literature and in earlier versions of
this work, arise when a held-out text is dated under a prior centred on the
scholarly date it is asked to recover.  For the most securely anchored
texts in the present corpus the linguistic evidence contributes under 6\%
of posterior precision, and the apparent recovery is a restatement of the
prior.  Corrected, the same design yields mean absolute errors near 100~yr
in Hebrew and near 120~yr in Greek.  We think it likely that this failure
mode is not confined to the present work, and we suggest that reported
holdout accuracies in this area be accompanied by the likelihood-only
estimate and by the fraction of posterior precision attributable to the
data.

What the corrected framework supports is narrower than what was previously
claimed and, we think, more useful.  It does not date texts: leave-one-out
error of 121--156~yr must be read against 137~yr for the trivial predictor
that assigns every text the corpus mean.  It does not assign texts to
historical periods with useful reliability, and it cannot resolve the
exilic period at all.  It does order texts, and it can determine on which
side of a well-separated boundary a text is likely to fall.  Applied to the
Pentateuchal sources under leakage-free conditions, that is enough to
establish that the Priestly, Deuteronomic and JE composites are post-exilic
with probability at least 0.92 under both model families and under a
correction for the genre imbalance of the training corpus --- and not
enough to place any of them within a particular post-exilic century.  For a
question that has been argued for a century and a half largely on
non-quantitative grounds, a bounded probabilistic answer to the coarse
question seems more valuable than an unbounded point estimate for the fine
one.
